%%
%% winter-lecture26.tex
%% 
%% Made by Alex Nelson <pqnelson@gmail.com>
%% Login   <alex@lisp>
%% 
%% Started on  2026-03-05T09:01:37-0800
%% Last update 2026-03-05T09:01:37-0800
%% 

\lecture{}

\begin{remark}
Every element of $K(B)$ can be represented by
$(E,\TrivialBundle{\CC}^{N})$ with no restriction on rank of $E$.
\end{remark}

\begin{node}
Elements $(E,\TrivialBundle{\CC}^{N})$ and
$(E',\TrivialBundle{\CC}^{N})$ in $\widetilde{K}(B)$ are the same if
there exists $k\in\NN$ and $\ell\in\NN$ such that
\begin{equation}\label{eq:math151b:lecture26:equivalence-relation-for-reduced-k-theory}
E\oplus\TrivialBundle{\CC}^{k}\iso E'\oplus\TrivialBundle{\CC}^{\ell}.
\end{equation}
\end{node}

\begin{node}
We could view $\widetilde{K}(B)$ as equivalence classes of vector
bundles over $B$ using Eq~\eqref{eq:math151b:lecture26:equivalence-relation-for-reduced-k-theory}
as the equivalence relation.
\end{node}

\begin{claim}
There is a bijection $\widetilde{K}(B)\stackrel{1:1}{\longleftrightarrow}[B,BU]$
where $U=U(\infty)$ and $BU$ is the classifying space of $U$.
\end{claim}

\begin{recall}
There is a ``universal'' $U$-bundle $U\into EU\onto BU$ such that $EU$
is contractible and moreover it is ``universal'' in the sense that: for
all compact Hausdorff spaces $X$ there is a bijection
\begin{equation*}
\begin{split}
\left\{\begin{matrix}\mbox{isomorphism classes of}
\mbox{$U$-bundles on }X
\end{matrix}\right\}&\stackrel{1:1}{\longleftrightarrow} [X,BU]\\
f^{*}(EU) & \mapsfrom [f]
\end{split}
\end{equation*}
given by sending the homotopy class $[f]\in[X,BU]$ to the pullback
bundle $f^{*}EU$ using any representating $f\in[f]$ (since the
pullback bundle is unique under homotopy, this works).
\end{recall}

\begin{node}
To prove the claim, we will need a notion of a frame bundle.
\end{node}

\begin{definition}
Given a $\CC^{n}$-vector bundle $\pi\colon E\to B$, the
\define{Associated Frame Bundle} $\AssocFrameBundle(E)$ is defined as
\begin{equation}
\AssocFrameBundle(E)=\{(b,v_{1},\dots,v_{n})\in
B\times\underbrace{E\times\cdots\times E}_{n\text{ times}}\mid
\{v_{1},\dots,v_{n}\}\mbox{ is a basis for fiber }
\pi^{-1}(\{b\})\subset E\}
\end{equation}
with the projection map $p\colon\AssocFrameBundle(E)\to B$ given by
sending $(b,v_{1},\dots,v_{n})\mapsto b$ to its first component.
\end{definition}

\begin{remark}
The associated frame bundle is a fiber bundle with fiber
$\GL_{n}(\CC)$ since any basis can be expressed as an element of
$\GL_{n}(\CC)$ applied to the canonical basis. (This is actually a
``principal $G=\GL_{n}(\CC)$-bundle'', i.e., there is a free and
transitive action of $G$ ($=\GL_{n}(\CC)$) on each fiber $\GL_{n}(\CC)$
by change of basis formula.) You can't fix an identity element on each
fiber---that would be the same data as a preferred basis.
\end{remark}

\begin{construction}
Given a fiber bundle $\pi\colon X\to B$ on $B$ with fiber $\GL_{n}(\CC)$ [and the
action of $\GL_{n}(\CC)$ on each fiber], you can build a
$\CC^{n}$-vector bundle in the following way: define the total space
$E$ as
\begin{subequations}
  \begin{align}
E &= X\otimes_{G}\CC^{n}\quad\mbox{where }G=\GL_{n}(\CC)\\
&=\{(x,v)\in X\times\CC^{n}\}/(Ax,v)\sim(x,Av)\mbox{ for }A\in\GL_{n}(\CC).
  \end{align}
\end{subequations}
Then $p\colon E\to B$ is given by $(x,v)\mapsto\pi(x)$. In this case,
we see $\AssocFrameBundle(E)\iso X$.
\end{construction}

\begin{remark}
Note: in this construction, you could use other groups $G$ instead of
$\GL_{n}(\CC)$, if you have an action $\rho\colon G\to\CC^{n}$ and you
then use $(Ax,v)\sim(x,\rho(A)v)$.
\end{remark}

\begin{remark}
This construction gives a bijection between isomorphism classes of
$\CC^{n}$-bundles over $B$ and isomorphism classes of $\GL_{n}(\CC)$
fiber bundles over $X$.
\end{remark}

\begin{fact}
If $E$ is a Hermitian vector bundle, then using the orthonormal basis
with respect to the Hermitian inner product on $E$ (in the
construction of the associated frame bundle, instead of ``just any old
basis'') we will get a $\U(n)$-bundle instead of $\GL_{n}(\CC)$-bundle.
\end{fact}

\begin{recall}
Now, we are in a position to prove the claim; what claim? Remember: There is a bijection $\widetilde{K}(B)\stackrel{1:1}{\longleftrightarrow}[B,BU]$
where $U=U(\infty)$ and $BU$ is the classifying space of $U$.
\end{recall}

\begin{proof}[Sketchy proof sketch]
We have equivalence relation
$E\sim E'$ if and only if there are $k,\ell\in\NN$ such that
$E\oplus\TrivialBundle{\CC}^{k}\iso E'\oplus\TrivialBundle{\CC}^{\ell}$.
We can think of this as a stabilization relation, where we send the
associated frame bundle's fiber $\U(m)\into\U(\infty)$ for $E$ and
similarly $\U(n)\into\U(\infty)$ for $E'$.

For any $\CC^{n}$-bundle $E$, you get a fiber bundle $U$ by the
inverse limit
\begin{equation}
\AssocFrameBundle^{\infty}(E):=\varprojlim(\AssocFrameBundle(E)\into\AssocFrameBundle(E)\oplus\TrivialBundle{\CC}^{1}\into\AssocFrameBundle(E)\oplus\TrivialBundle{\CC}^{2}\into\AssocFrameBundle(E)\oplus\TrivialBundle{\CC}^{3}\into\dots)
\end{equation}
Then $E\sim E'$ if and only if $\AssocFrameBundle^{\infty}(E)\iso\AssocFrameBundle^{\infty}(E')$
as $\U(\infty)$-fiber bundles. Therefore, we have the bijection
\begin{equation}
\widetilde{K}(B)\stackrel{1:1}{\longleftrightarrow}\left\{\begin{matrix}\mbox{isomorphism classes of}\\ 
U(\infty)\mbox{-bundles over } B\end{matrix}\right\}\stackrel{1:1}{\longleftrightarrow}[B,BU].
\end{equation}
Hence the sketchy proof sketch.
\end{proof}

\begin{remark}
The unreduced K-theory version has the bijective correspondence
\begin{equation*}
K(B)\stackrel{1:1}{\longleftrightarrow}[B,\ZZ\times BU].
\end{equation*}
\end{remark}

\begin{example}
The K-theory for a point $K(\point{*})\iso\ZZ$ since
$(E,E')\mapsto\dim(E)-\dim(E')$ is the only thing we can really do here.
\end{example}

\begin{example}
A zero-sphere is just two points, so $K(\sphere{0})\iso\ZZ\oplus\ZZ$.
\end{example}

\begin{example}
For a 1-sphere, $K(\sphere{1})\iso\ZZ$. The reasoning is as follows:

\textsc{Claim:} All complex vector bundles over $\sphere{1}$ are
trivial (in the strict vector bundle sense).

\begin{proof}
It's because $\GL_{n}(\CC)$ is connected, so $\GL_{n}(\CC)\homotopic\U(n)$,
in contrast to $\GL_{n}(\RR)$ which has two connected components (for
orientation-reversing and orientation-preserving transformations,
since $\det\colon\GL_{n}(\RR)\to\RR\setminus\{0\}$). Then since
$\sphere{1}=I_{1}\cup I_{2}=I/0\sim 1$ gives the transition map at
$0\in I$ to an element of $\GL_{n}(\CC)$ which can be homotoped to the
identity matrix. This implies the $\CC^{n}$-bundle over $\sphere{1}$
is trivial.
\end{proof}
\end{example}

\begin{example}
We see $K(\sphere{2})\iso\ZZ\oplus\ZZ\iso\widetilde{K}(\sphere{2})\oplus\ZZ$
where $\widetilde{K}(\sphere{2})\stackrel{1:1}{\longleftrightarrow}[\sphere{2},BU]$.
But
\begin{subequations}
  \begin{align}
[\sphere{2},BU]
&\iso[\Sigma\sphere{1},BU]\\
&\iso[\sphere{1},\Omega BU]\mbox{ by adjoint property}\\
&\iso[\sphere{1},\U(\infty)] \mbox{ by long exact sequence of homotopy
  groups for fibration}\\
&\iso\pi_{1}(\U(\infty)).
  \end{align}
\end{subequations}
Then for the inclusion $\U(n)\into\U(n+1)$, there is a fiber bundle
$\U(n)\into\U(n+1)\onto\sphere{2n+1}$ related to picking a basis in
$\CC^{n}$. Then $i_{*}\colon\pi_{1}(\U(n))\into\pi_{1}(\U(n+1))$ is an
isomorphism for all $n$, and since
\begin{equation}
\pi_{1}(\U(\infty))\iso\pi_{1}(\U(1))\iso\pi_{1}(\sphere{1})\iso\ZZ,
\end{equation}
we conclude $\widetilde{K}(\sphere{2})\iso\ZZ$.

What is the generator for it? Well, we can find it by the following:
recall the Clutching construction, we have the $\CC$-vector bundles
over $\sphere{2}$ is in bijective correspondence with elements of $[\sphere{1},\GL_{1}(\CC)\iso\sphere{1}]$,
and the generator of this is given by just the identity map on the
circle.
This equivalence class $[\id]\in[\sphere{1},\GL_{1}(\CC)]$ corresponds
to the $\CC^{1}$-bundle denoted $H$, which corresponds to the
tautological line bundle for $\sphere{2}\iso\CP^{1}$.

(\textbf{Exercise:} prove $(H\otimes
H)\oplus\TrivialBundle{\CC}^{1}\iso H\oplus H$. [Hint: look at
  transition map])

What is the ring structure of $K(\sphere{2})$? We see that this is
given by
\begin{equation}
K(\sphere{2})\iso\ZZ[H]/(H^{2}-1)\quad\mbox{as rings}.
\end{equation}
We also have the isomorphism \emph{as groups}
\begin{equation}
K(\sphere{2})\iso\ZZ\oplus\ZZ\quad\mbox{as groups}
\end{equation}
where one $\ZZ$ is generated by $H$, and the other $\ZZ$ is generated
by $1$ (the multiplicative unit).
\end{example}

\begin{proposition}
Let $X$ be a compact Hausdorff space and $A\subset X$ is a closed
subspace with the inclusion map $i\colon A\into X$ and the quotient
map $q\colon X\onto X/A$. Then this gives us a sequence
\begin{equation}
\widetilde{K}(X/A)\xrightarrow{q^{*}}\widetilde{K}(X)\xrightarrow{i^{*}}\widetilde{K}(A),
\end{equation}
which is exact at $\widetilde{K}(X)$.
\end{proposition}

\begin{remark}
This is the key step for proving we have a long exact sequence for K-theory.
\end{remark}

\begin{remark}
Here the intuition is that $q^{*}$ ``fills'' $A$ with trivial bundle,
and $i^{*}$ restricts attention to $A$.
\end{remark}